\documentclass[sigconf,nonacm]{acmart}

%% Axiomander -- Vericoding paper (FSE-targeted draft)
%% ACM acmart class, sigconf (two-column) format.
%% nonacm: this is a self-archived draft, not a submitted ACM proceedings copy.

\usepackage{booktabs}
\usepackage{mathtools}
\usepackage{mathpartir}
\usepackage{listings}
\usepackage{xcolor}
\let\Bbbk\relax
\usepackage{amssymb}

%% Notation
\newcommand{\sepc}{\mathbin{\ast}}
\newcommand{\wpre}{\mathrm{wp}}
\newcommand{\hoare}[3]{\{#1\}\ #2\ \{#3\}}
\newcommand{\code}[1]{\texttt{#1}}

%% Lightweight code listing style for Python / Coq snippets
\lstdefinestyle{snippet}{
  basicstyle=\ttfamily\footnotesize,
  breaklines=true,
  columns=fullflexible,
  keepspaces=true,
  showstringspaces=false,
  frame=single,
  framesep=4pt,
  rulecolor=\color{gray!40},
  xleftmargin=2pt,
  xrightmargin=2pt,
}
\lstset{style=snippet}

\begin{document}

\title{Vericoding: Specification-Centric Software Development for the AI Era}
\subtitle{How cheap implementation makes specifications the primary artifact}

\author{Gavin Mendel-Gleason}
\affiliation{%
  \institution{Scidonia}
  \country{}
}
\email{}

\renewcommand{\shortauthors}{Mendel-Gleason}

\begin{abstract}
Large language models have sharply reduced the cost of producing implementations
but have not reduced the cost of deciding whether an implementation satisfies
intent. Mainstream software engineering still treats source code as the primary
maintained artifact, with specifications, tests, and proofs as secondary
validation. We argue that the changed economics invert this relationship. When
implementation is cheap and verification stays expensive, the scarce resource is
\emph{specification}, and the durable artifact should be the specification rather
than the code. We call the resulting model \emph{vericoding}: developers
iteratively refine a strong specification, an AI system synthesizes the
implementation and the subsidiary specifications its proof needs, and a verifier
returns a deterministic verdict -- with a structured residual, not a vague
complaint, on failure. We make the case for this model, identify the properties
it demands of a verifier (compositional reasoning so that proofs behave like
cached build artifacts, and structured failure so that an AI can repair against a
proof obligation rather than re-guess against a prompt), and describe its
realization in Axiomander, a working verifier for Python whose contracts are
ordinary \code{assert} statements. We give an evaluation methodology built around
retrofitting, synthesis, and maintenance, report early evidence from the
prototype, and discuss the human-factors and soundness questions the model
raises. The central claim is economic rather than technical: software development
in the AI era is better understood as specification management than as
programming.
\end{abstract}

\keywords{vericoding, specification-centric development, deductive verification,
LLM code synthesis, contracts, compositional verification, separation logic,
software engineering process}

\maketitle

\section{Introduction}
\label{sec:intro}

For most of the history of software engineering, writing code was the expensive
step. Processes, tools, and review practices grew up around that assumption: the
source code is the artifact a team reads, edits, owns, and trusts, and tests
sample its behaviour from the side. Specifications, where they exist at all, are
usually informal.

Large language models change the cost structure. Generating a plausible
implementation of a described task is now fast and cheap. Deciding whether that
implementation is \emph{correct} is not: it is exactly as hard as it ever was,
and arguably harder, because an LLM is trained to satisfy the \emph{prompt}, not
the \emph{intent}. Its characteristic failure is silent -- code that compiles,
runs on the example, and reads correctly to a skimming reviewer, while quietly
mishandling the inputs nobody thought about. The bottleneck has moved from
implementation effort to intent specification.

This paper takes that shift seriously and asks what development looks like when
specification, not code, is the scarce and therefore primary artifact. We call
the model \emph{vericoding}:

\begin{center}
\emph{Traditional:}\quad Human $\rightarrow$ Code $\rightarrow$ Tests $\rightarrow$ Confidence\\[2pt]
\emph{Vericoding:}\quad Human $\rightarrow$ Specification $\rightarrow$ AI $\rightarrow$ Implementation $\rightarrow$ Verification
\end{center}

The human writes and maintains a strong specification. An AI derives an
implementation -- and, freely, the helper functions and \emph{their}
specifications that its proof needs. A verifier decides, deterministically,
whether the code meets the specification, and on failure hands back a structured
residual that drives the next attempt. The implementation becomes a derived,
disposable artifact; the specification persists across many implementations.

Two properties make this loop tractable rather than aspirational, and they are
the engineering content of the paper. First, \textbf{compositional verification}:
a caller must be provable from a callee's contract, not its body, so that editing
a body does not re-verify the world. Under this discipline a proof behaves like a
cached object file and the verifier behaves like a build system for correctness.
Second, \textbf{structured failure}: a negative verdict must return the exact
remaining obligation (with hypotheses, and a counterexample where one exists), so
that the AI repairs against a goal rather than re-guessing against a prompt.

\paragraph{Contributions.}
\begin{itemize}
  \item We articulate vericoding as a software-engineering model and state its
  four operating principles (Section~\ref{sec:model}), framing the shift as an
  economic consequence of cheap implementation.
  \item We argue that the model stands or falls on \emph{compositional scaling},
  and show how contracts-as-interfaces plus proof caching turn verification into
  an incremental build system for correctness (Section~\ref{sec:compose}).
  \item We describe the \emph{heal loop} -- residual-driven repair -- and why it
  dominates naive re-prompting (Section~\ref{sec:heal}).
  \item We report a concrete realization, Axiomander, a verifier for Python whose
  contracts are plain \code{assert} statements and whose verdict bottoms out in a
  machine-checked proof (Section~\ref{sec:axiomander}), and give an evaluation
  methodology with early evidence (Section~\ref{sec:eval}).
  \item We surface the human-factors and soundness questions the model raises
  (Sections~\ref{sec:human}--\ref{sec:limits}).
\end{itemize}

\section{The traditional model and its limits}
\label{sec:traditional}

Traditional development is code-centric: the human writes code, tests sample it,
and confidence accumulates through review. This made sense when humans were the
code generators -- the artifact you paid to produce was the artifact you kept.
The weakness is in what the validation mechanisms actually guarantee. Types
constrain the \emph{shape} of data but say nothing about the relationship between
inputs and outputs; tests assert facts at finitely many points. Both are
\emph{samples} of intent, not statements of it.

\begin{table}[t]
\caption{A ladder of guarantees. Each mechanism catches more, but only a strong
specification states intent totally.}
\label{tab:ladder}
\small
\begin{tabular}{@{}lll@{}}
\toprule
Mechanism & Catches & Misses \\
\midrule
Prompt review & gross misunderstanding & everything subtle \\
Types & shape errors & values, relations, invariants \\
Tests & cases you thought of & cases you did not \\
\textbf{Strong spec.} & anything expressible & only what you leave unstated \\
\bottomrule
\end{tabular}
\end{table}

LLM-generated code stresses exactly this gap (Table~\ref{tab:ladder}). The model
is not adversarial; it is \emph{optimistic}. It assumes the inputs it did not
consider behave like the ones it did, that an unwritten invariant holds, and that
an omitted error path is unreachable. A green test suite says ``these inputs are
right''; it cannot say ``every input is right.'' As implementation volume rises
and human reading time does not, sampling-based validation falls further behind.
What is needed is a mechanism that decides correctness across \emph{all} inputs,
deterministically -- which is what a proved specification provides.

\section{The vericoding model}
\label{sec:model}

Vericoding inverts the data-flow of development. Rather than writing code and
sampling it, the developer writes a specification and derives code beneath it:

\[
\underbrace{(P, Q)}_{\text{spec, human-authored}} \;\leadsto\;
\underbrace{c}_{\text{code, AI-derived}}
\quad\text{such that}\quad \vDash \hoare{P}{c}{Q}.
\]

The human authors both halves: the assumption $P$ on inputs and the guarantee $Q$
on outputs. The code is derived so that $Q$ holds \emph{under} $P$. Four
principles follow.

\paragraph{Principle 1: specification is the control surface.} You steer the
system by editing specifications, not code. The reviewable surface shrinks from
``all the code'' to ``the contracts at the boundary.''

\paragraph{Principle 2: code is a derived artifact.} The implementation is
disposable. If it does not verify, regenerate it. We do not trust the code
because we read it; we trust it because it was \emph{proved} against a
specification we did read.

\paragraph{Principle 3: correctness requires proof.} A specification is only as
good as our ability to discharge it. Runtime contracts catch violations after
deployment on whatever inputs occur; static proof catches them before deployment
on all inputs. The verdict is binary and deterministic.

\paragraph{Principle 4: failure returns structured information.} A verifier that
says only ``no'' is useless to an automated repair loop. The verdict must carry
the residual obligation and, where the goal is false, a concrete counterexample.

These principles only pay off if two scaling problems are solved -- composition
(Section~\ref{sec:compose}) and structured repair (Section~\ref{sec:heal}) -- and
if a real verifier can implement them over a mainstream language
(Section~\ref{sec:axiomander}).

\section{Specification management}
\label{sec:specmgmt}

If code is derived, the central human activity is no longer editing
implementations but \emph{managing specifications}: refining contracts, reviewing
them, decomposing them, and evolving requirements. This shifts effort toward
artifacts that are smaller, more declarative, and easier to review than the code
they constrain.

A key degree of freedom belongs to the AI, not the human. To prove a top-level
entry point, the AI may decompose the problem into helper functions, each with
its \emph{own} contract; those sub-specifications are how it makes the global
proof go through, and they are checked the same way -- there is no privileged,
unverified layer. The human owns the boundary contract; the AI owns the interior
decomposition. Formally, a caller's proof uses only a callee's contract:
\[
\inferrule{\hoare{P_g}{g}{Q_g} \\\\ \hoare{P_f}{f}{Q_f}\ \text{using}\ Q_g}
{\hoare{P_f}{f}{Q_f}}
\]
The body of $g$ never appears in $f$'s proof. This is what makes the human's
review surface small even when the generated interior is large, and it is the
same property that makes the verifier incremental (Section~\ref{sec:compose}).

\section{Realization: Axiomander}
\label{sec:axiomander}

Vericoding is a methodology, not a tool, but a methodology is only credible if it
runs. Axiomander is a working realization for Python. Its trust kernel is the
Rocq proof assistant (formerly Coq); separation logic, via the Iris framework,
supplies the weakest-precondition calculus and the heap reasoning that Python's
mutation and aliasing demand; and an LLM acts as an oracle \emph{inside} the
loop, never as the trust kernel. The methodology is backend-agnostic -- any
engine that returns a sound, deterministic decision plus a structured residual
would serve -- but the realization commits to Rocq{+}Iris because separation
logic is the right foundation for resource ownership in real code.

\paragraph{Contracts are ordinary Python.} The user writes no imports and no
decorators. Contracts are plain \code{assert} statements, classified by position,
or a verifier-only \code{axiomander:} docstring:

\begin{lstlisting}[language=Python]
def clamp(val: int, lo: int, hi: int) -> int:
    assert lo <= hi                      # precondition
    if val < lo:    result = lo
    elif val > hi:  result = hi
    else:           result = val
    assert lo <= result <= hi            # postcondition
    assert implies(val < lo, result == lo)
    assert implies(val > hi, result == hi)
    return result
\end{lstlisting}

Leading assertions are preconditions, trailing assertions after the result are
postconditions, and loop-body assertions are invariants. A shared linter lowers
both syntaxes to one contract IR, so the user's code stays dependency-free while
the verifier does the work.

\paragraph{From Python to a verdict.} The pipeline lowers the contract IR to an
intermediate imperative language, emits a weakest-precondition obligation, and
discharges it through a ladder that runs cheap, deterministic tactics first and
escalates only on residuals (Table~\ref{tab:pipeline}). The defining theorem --
that a Hoare triple holds iff the precondition implies the weakest precondition
of the body -- is proved sound in Rocq/Iris, so an accepted verdict is
trustworthy by construction, and the LLM, which may be wrong, only ever proposes
proofs that Rocq re-checks.

\begin{table}[t]
\caption{Axiomander's proof pipeline. Deterministic tactics run first; the LLM is
the fallback, not the first resort.}
\label{tab:pipeline}
\small
\begin{tabular}{@{}cll@{}}
\toprule
Tier & Mechanism & Handles \\
\midrule
1  & deterministic Rocq (\code{wp\_reduce}, \code{lia}) & assignments, branches, loops, frames \\
2  & SMT / Rocq-hammer & linear/nonlinear arithmetic, FOL \\
2b & theory-SMT (\code{QF\_SLIA}) & strings, regex, float dimensions \\
3  & LLM oracle & invariants, induction, lemma chains \\
\bottomrule
\end{tabular}
\end{table}

\section{Compositional scaling: a build system for correctness}
\label{sec:compose}

Vericoding lives or dies on composition. If a one-line spec edit re-verified the
whole program, the loop would not scale to real codebases. The discipline,
borrowed from incremental build systems, is:
\begin{center}
\emph{Body changes invalidate local proofs; contract changes invalidate callers.}
\end{center}

A contract is an interface; a proof is a cached build artifact. Each function
exports a semantic summary
\[
\mathrm{summary}(g) = \mathrm{hash}\big(\text{pre},\ \text{post},\ \text{reads},\ \text{writes},\ \text{raises}\big),
\]
and a caller $f$ depends on $\mathrm{summary}(g)$, not on $g$'s implementation. If
$g$'s body changes but its summary is stable, $f$'s proof is reused; if the
summary changes, only $f$ and its transitive callers are re-verified
(Table~\ref{tab:incremental}). The verifier becomes an incremental build system
for correctness, not a batch theorem prover, and the same opaque-contract
boundary serves three purposes at once: modular verification, proof reuse under
iteration, and \emph{stubbing} -- an unverified library is replaced by its
declared contract, which becomes an axiom available to callers, with the trust
boundary made explicit and reviewed once.

\begin{table}[t]
\caption{Incremental invalidation. The unit of reuse is the contract, not the
file.}
\label{tab:incremental}
\small
\begin{tabular}{@{}ll@{}}
\toprule
Change & Re-verify \\
\midrule
body only & the function \\
local assert / invariant & the function \\
\textbf{contract} & the function $+$ transitive callers \\
\textbf{callee contract} (summary hash) & the caller \\
\bottomrule
\end{tabular}
\end{table}

\paragraph{Composition needs the right logic.} The frame property -- reasoning
about a component is insensitive to state it does not touch -- is what makes
proofs local. Whether it is cheap depends on the assertion logic. In a companion
study we built the same backend twice: once over a flat name-to-value store,
where locality has to be \emph{reconstructed} by hand, and once over Iris
separation logic, where it is the native shape of the logic. The difference was
stark and is worth stating concretely, because it is the difference between
vericoding scaling and not. A two-call function (two increments and a sum)
required, over the flat store, two explicit frame applications plus fourteen
variable-preservation conjuncts threaded through intermediate postconditions and
a 29-line, hand-decomposed proof that the monolithic tactic could not close; the
separation-logic backend discharged the analogous chain with \emph{zero} framing
obligations in an 11-line proof that the generator emitted and the kernel
re-checked. The gap is not only efficiency. Functions that mutate an owned field,
or mutate state and then raise through it, have \emph{no} flat-store rendering at
all -- the flat store cannot express ownership of a mutable cell. The
substructural logic is what makes such programs \emph{verifiable in the first
place}, and realistic Python is full of them. Composition, in other words, is
both a scaling property and an expressivity property.

\section{The heal loop}
\label{sec:heal}

Naive AI coding is a guessing loop: prompt, code, ``try again,'' code. The model
gets no signal it can act on. Vericoding replaces this with structured repair:
\begin{center}
spec $\rightarrow$ code $\rightarrow$ verify $\rightarrow$ residual goal $\rightarrow$ repair $\rightarrow\ \cdots$
\end{center}

Two things distinguish a heal loop from re-prompting. A \textbf{residual goal} is
the exact remaining obligation with its hypotheses; a \textbf{counterexample} is
a concrete, typed witness of failure. Consider a contract whose pre- and
post-conditions are regular-expression membership:

\begin{lstlisting}[language=Python]
def accept(p: str) -> str:
    """ axiomander:
        requires: p.re_match("[0-9]{3}")
        ensures:  result.re_match("[a-z]+") """
    result = p
    return result
\end{lstlisting}

The verifier does not say ``wrong.'' It returns a witness -- \code{result =
"000"} satisfies the requires-regex and violates the ensures-regex -- because a
digit string can never match \code{[a-z]+}. The AI repairs against that fact.
Crucially, obligations are kept small and stably identified (e.g.
\code{sort.insert.preserves\_sorted.branch\_2}) rather than fused into one giant
theorem, so a failed stage yields a reusable artifact -- a named residual with
its hypotheses -- rather than a dead end, and per-obligation caching,
parallelism, and narrow AI prompts all follow from the same structure.

\section{Evaluation methodology}
\label{sec:eval}

Vericoding makes empirical claims -- that the maintained surface shrinks, that
iteration stays local, that structured repair converges -- and they must be
measured rather than asserted. We propose three case studies and the metrics
each isolates, then report what the current prototype already demonstrates.

\paragraph{Case study 1: retrofitted verification.} Add contracts to an existing
Python codebase and measure contracts written per function, proof obligations
generated, and the fraction discharged at each tier. This isolates the cost of
\emph{adopting} the model on code not written for it.

\paragraph{Case study 2: synthesis from specification.} Generate implementations
from specifications alone and measure repair iterations to a verified result,
per-tier proof-success rates, and residual complexity at the point the AI takes
over. This isolates the heal loop's convergence.

\paragraph{Case study 3: maintenance under changing requirements.} Take a
requirement change and run it both ways -- traditional (edit code, update tests)
and vericoding (edit specification, regenerate, re-verify) -- measuring
engineering effort and the blast radius of re-verification. This isolates the
incremental-build claim of Section~\ref{sec:compose}.

\paragraph{Early evidence.} The prototype already exercises the mechanisms these
studies will quantify, which is what makes the studies runnable rather than
hypothetical. Axiomander verifies end-to-end on a corpus of contracted functions;
on that corpus the deterministic tier discharges the large majority of goals
(roughly four-fifths in our runs) before any external prover is invoked, with SMT
and the LLM handling the residual. It returns typed counterexamples for false
goals, including the regex case above, and it dogfoods: more than a dozen of
Axiomander's own functions -- covering enum dispatch, \code{isinstance}-based
case analysis, string methods, and object construction -- carry contracts and are
verified by the tool itself. The compositional measurement of
Section~\ref{sec:compose} is a controlled comparison on the same functions across
two backends. We are explicit that this is prototype-stage evidence, not a
completed controlled study; the three case studies above, on third-party code and
with human-effort measurement, are the work that turns it into one.

\section{Human factors}
\label{sec:human}

Vericoding makes a falsifiable bet about where engineers spend attention. The
hypothesis is that traditional development asks humans to read large volumes of
implementation, whereas vericoding asks them to read comparatively small contract
surfaces, and that the durable artifacts (specifications) and the churning
artifacts (implementations) separate cleanly. The questions worth measuring are
direct: what do engineers actually review, which artifacts stay stable across a
feature's life, which churn, and how much generated code a human must read before
trusting a verified result. If, in practice, reviewers feel compelled to read the
generated code anyway, the central economic claim weakens; if they review
contracts and trust proofs, it holds. These are usability and trust questions, and
they are as central to the model's success as the logic underneath it.

\section{Limitations and threats to validity}
\label{sec:limits}

Vericoding shifts effort; it does not eliminate it. The most important shift is
also the most important risk: the specification is now the artifact that can be
wrong. A proof against an incorrect specification certifies the wrong thing, and a
specification can be vacuous (an unsatisfiable precondition proves anything).
Mitigations -- sanity-checking preconditions for satisfiability, surfacing
vacuity, and reviewing contracts as the primary artifact -- reduce but do not
remove this risk.

Three further limitations bound the model. \emph{Soundness of assumptions}:
SMT-discharged facts imported as axioms, and library stubs taken on faith, are
trust boundaries; they must be logged and auditable, and a wrong stub is a silent
hole. \emph{Coverage}: a verifier supports a language fragment, and constructs
outside it (runtime metaprogramming, reflection) must be gated rather than
silently mistranslated. \emph{Scalability of proof search}: the LLM tier can fail
to find an invariant or a lemma chain, in which case the loop stalls and a human
must intervene. Finally, our evaluation is at present prototype-stage and
self-referential; the controlled case studies of Section~\ref{sec:eval}, on
code we did not write, are required before the model's effort claims can be
considered established.

\section{Related work}
\label{sec:related}

Deductive verification has a long lineage -- Hoare logic~\cite{hoare1969},
Dijkstra's weakest preconditions~\cite{dijkstra1975}, and separation
logic~\cite{reynolds2002} as realized in modern frameworks such as
Iris~\cite{jung2018iris}. Auto-active verifiers including
Dafny~\cite{leino2010dafny}, F$^\star$~\cite{swamy2016fstar}, and
Verus~\cite{lattuada2023verus} let programmers discharge functional-correctness
obligations with substantial automation. Axiomander reuses this machinery; its
departure is not in the logic but in the \emph{process context}. Where these
systems expose the assertion logic to the programmer and treat the annotated
program as the artifact, vericoding keeps the logic an implementation detail
(contracts are plain Python), makes the specification the maintained artifact,
and places an LLM in a structured repair loop over residual proof state.

On the synthesis side, large language models generate code at scale~\cite{chen2021codex}
but optimize for prompt resemblance, not verified intent; the gap they leave is
precisely the one a verifier closes. The contribution we claim is a coherent
software-engineering model that joins these threads: cheap AI synthesis under a
human-maintained specification, made to scale by compositional, contract-keyed
proof caching, and made convergent by residual-driven repair.

\section{Conclusion}
\label{sec:conclusion}

The observation behind vericoding is economic, not technical. When implementation
becomes inexpensive and verification stays expensive, the specification becomes
the scarce resource and, naturally, the primary maintained artifact. The
engineering consequences are concrete: the reviewable surface shrinks to
contracts, proofs become cached build artifacts invalidated along contract
boundaries, and AI repair is driven by proof obligations rather than prompts. A
working verifier for Python, with contracts written as ordinary assertions, shows
the model is implementable today. The open question we find most interesting is
the one the model poses by construction: whether software development in the AI
era is best understood as programming -- or as specification management.

\begin{thebibliography}{9}

\bibitem{hoare1969}
C.~A.~R. Hoare.
\newblock An axiomatic basis for computer programming.
\newblock \emph{Communications of the ACM}, 12(10):576--580, 1969.

\bibitem{dijkstra1975}
E.~W. Dijkstra.
\newblock Guarded commands, nondeterminacy and formal derivation of programs.
\newblock \emph{Communications of the ACM}, 18(8):453--457, 1975.

\bibitem{reynolds2002}
J.~C. Reynolds.
\newblock Separation logic: A logic for shared mutable data structures.
\newblock In \emph{Proc.\ 17th IEEE Symposium on Logic in Computer Science
(LICS)}, pages 55--74, 2002.

\bibitem{jung2018iris}
R.~Jung, R.~Krebbers, J.-H. Jourdan, A.~Bizjak, L.~Birkedal, and D.~Dreyer.
\newblock Iris from the ground up: A modular foundation for higher-order
concurrent separation logic.
\newblock \emph{Journal of Functional Programming}, 28:e20, 2018.

\bibitem{leino2010dafny}
K.~R.~M. Leino.
\newblock Dafny: An automatic program verifier for functional correctness.
\newblock In \emph{Proc.\ 16th Int.\ Conf.\ on Logic for Programming, Artificial
Intelligence and Reasoning (LPAR-16)}, pages 348--370, 2010.

\bibitem{swamy2016fstar}
N.~Swamy, C.~Hri\c{t}cu, C.~Keller, A.~Rastogi, A.~Delignat-Lavaud, et~al.
\newblock Dependent types and multi-monadic effects in F$^\star$.
\newblock In \emph{Proc.\ 43rd ACM SIGPLAN-SIGACT Symposium on Principles of
Programming Languages (POPL)}, pages 256--270, 2016.

\bibitem{lattuada2023verus}
A.~Lattuada, T.~Hance, C.~Cho, M.~Brun, I.~Subasinghe, Y.~Zhou, J.~Howell,
B.~Parno, and C.~Hawblitzel.
\newblock Verus: Verifying Rust programs using linear ghost types.
\newblock \emph{Proc.\ ACM on Programming Languages}, 7(OOPSLA1), 2023.

\bibitem{chen2021codex}
M.~Chen, J.~Tworek, H.~Jun, Q.~Yuan, H.~P.~de~Oliveira~Pinto, et~al.
\newblock Evaluating large language models trained on code.
\newblock \emph{arXiv preprint arXiv:2107.03374}, 2021.

\end{thebibliography}

\end{document}
